% 依据用户于2026-11-05提供的勘误说明编写；编号与该说明一一对应。
% 本文件由主文件在导言区载入。正文保留原式，在相应位置调用译者补注。
\DeclareTranslatorNote{1}{Picard 下降与可分闭包}{这里的下降正合列应采用可分闭包 $k^s$，即
\[
0\longrightarrow\Pic(X)\longrightarrow\Pic(X_{k^s})^{\Gal(k^s/k)}\longrightarrow\Br(k).
\]
若坚持使用代数闭包，则须增加 $k$ 完美的假设；$X(k)\ne\varnothing$ 时的同构断言也应作同样处理。例如，取 $k=\mathbb F_2(t)^{\mathrm{sep}}$ 和 $E:y^2+y=x^3$。此时 $k$ 可分闭，$G=1$ 且 $\Br(k)=0$，但 $\sqrt t\notin k$；在 $\bar k$ 上可取 $P=(\sqrt t,y)$，于是 $\mathcal O(P-O)$ 不能下降到 $k$。因此代数闭包上的 Picard 群不能由所写的不动部分下降式给出。命题 1.1 本身的拉回单射结论仍可保留。}

\DeclareTranslatorNote{2}{定理 6.4 的正特征证明缺口}{本定理的陈述包含正特征基域，但下列证明没有覆盖其全部范围：a) 所用的平展上同调单射，先前只在 $\ell\ne\operatorname{char}k$ 的初级挠元上建立；b) 对剩余核的指数有界论证又明确假设特征为零。因此，正特征情形，特别是 $p$ 初级挠元，不能由这里的证明推出。若仅采用本书给出的论证，应把 $k$ 限为 $\mathbb Q$ 上有限生成的域；保留原陈述则需另补正特征证明。这是证明范围的缺口，不能只靠改动一个符号补足。}

\DeclareTranslatorNote{3}{Brauer 群目标应取模常数类的商}{这里的自然对应应以商群为目标。在本段假设下，Hochschild--Serre 谱序列给出
\[
H^1(k,\Pic\bar X)\simeq\Br_1(X)/\im\Br(k),
\]
其中 $\Br_1(X)=\Ker[\Br(X)\to\Br(\bar X)]$。没有另行指定提升所需的结构，不能自然地把这一商群提升到 $\Br(X)$。因此所写复合应在上述商群上解释，相应交换图也应据此取商。常数 Brauer 类与次数为零的闭链群 $A_0(X)$ 的配对为零，故所需配对可自然下降到这个商群。}

\DeclareTranslatorNote{4}{定理 9.1 中的 $H^0/H^2$}{第二个交换图及紧接说明中的两处 $H^0(\bar C,\mathbb Q/\mathbb Z(2))^G$，均应读为 $H^2_{\et}(\bar C,\mathbb Q/\mathbb Z(2))^G$。这里用的是谱序列微分 $d_2^{0,2}:E_2^{0,2}\to E_2^{2,1}$，而曲线的迹同构给出
\[
H^2_{\et}(\bar C,\mathbb Q/\mathbb Z(2))\simeq\mathbb Q/\mathbb Z(1).
\]
真正的 $H^0$ 则是 $\mathbb Q/\mathbb Z(2)$，其 Tate 扭转次数不同，不能由此认作 $\mu(k)$。例如在 $k=\mathbb Q$ 时，错误表达式的不动部分含三阶元，而 $\mu(\mathbb Q)=\{\pm1\}$。}

\DeclareTranslatorNote{5}{$p$ 处 Euler 因子的 Frobenius}{在 $\ell=p$ 的分支，若使用第二章 §1.3 中标准的 Fontaine Frobenius，则应写
\[
\det_\Lambda(1-\varphi_p t;D_{\mathrm{crys}}(V)),
\]
不应把 $\varphi_p$ 取逆。例如 $V=\mathbb Q_p(1)$ 时，$D_{\mathrm{crys}}(V)$ 上的 Frobenius 为 $p^{-1}$，因子应为 $1-p^{-1}t$，现有写法却给出 $1-pt$。这一修正也与本节的 $L(V(r),s)=L(V,s+r)$ 相容。$\ell\ne p$ 分支中的 $\sigma_\ell^{-1}$ 保持原约定。}

\DeclareTranslatorNote{6}{$V$ 的定义缺少 Tate 扭转}{此处应补上 $(r)$：
\[
V=H^m_{\et}(X\otimes_{\mathbb Q}\overline{\mathbb Q},\mathbb Q_p)(r).
\]
紧接着的例子取 $X=\Spec K$、$(m,r)=(0,1)$，并使用 $\mathbb Q_p(1)$ 系数的上同调；原定义不带扭转，只会给出 $\mathbb Q_p$。例如 $K=\mathbb Q$ 时，两种系数的零次上同调分别为 $\mathbb Q_p$ 与 $0$。猜想 2.2.7 的定义已写出所需的 $(r)$。}

\DeclareTranslatorNote{7}{猜想 2.2.7 的例外参数}{条件 (ii) 中第一组例外应为 $(m,r)=(0,1)$，即改为
\[
(m,r)\ne(0,1),(2\dim X,\dim X+1).
\]
这也关系到 2.2.6(4) 中 $H^1_S\simeq H^1_{\mathrm{lim}}$ 的断言。按补注 6 补上 Tate 扭转后，取 $X=\mathbb P^1_{\mathbb Q}$、$m=0,r=1$，原写法仍允许该例，但
\[
H^1_S\simeq\mathbb Z[1/S]^\times\otimes\mathbb Q_p,
\qquad H^1_{\mathrm{lim}}\simeq\mathbb Q^\times\otimes\mathbb Q_p.
\]
前者有限维，后者无限维，不可能同构。因此这里必须排除 $(0,1)$。}

\DeclareTranslatorNote{8}{Lichtenbaum 公式中的 $S$ 因子}{右侧取 $\mathcal O_K[1/S]$ 上的上同调时，左侧应为去掉 $S$ 处 Euler 因子的 $\zeta_K^{(S)}(1-r)_{(p)}$，沿用前文对部分 zeta 函数的定义。若保留完整 $\zeta_K$ 的写法，则须相应限制右侧的 $S$，例如取 $S=\{p\}$。取 $K=\mathbb Q,p=3,r=2$，把 $7$ 加入 $S$ 会引入因子 $1-7=-6$，其 $3$ 幂部分为 $3$，定位正合列使右侧增加同样的因子，而完整 zeta 值不随 $S$ 改变。}

\DeclareTranslatorNote{9}{分圆单位的素数幂判据}{对 $N\ge2$，此处“当且仅当 $N$ 为合数”应为“当且仅当 $N$ 不是素数幂”。例如 $N=4$ 虽为合数，但 $c_4=(1-i)(1+i)=2$，在 $\mathbb Z[i]^+=\mathbb Z$ 中不是单位。正确判据由分圆多项式 $\Phi_N(1)$ 的取值立即得到；后一句关于在 $\mathbb Z[\alpha,1/N]^+$ 中为单位的叙述不受这一端点判据修正的影响。}

\DeclareTranslatorNote{10}{有限域 zeta 函数的零极点阶符号}{这里以零点的阶为正、极点的阶为负，则两处 $(-1)^i$ 都应改为 $(-1)^{i+1}$：
\[
\operatorname{ord}_{t=1}Z_\Lambda(X,D_\Lambda\mathcal F,t)
=\sum_i(-1)^{i+1}\rank_\Lambda P^i
=\sum_i(-1)^{i+1}\rank_\Lambda Q^i.
\]
例如 $X=\Spec\mathbb F_\ell$、$\Lambda=\mathbb Q_p$、$\mathcal F=\mathbb Q_p$ 时，$Z=(1-t)^{-1}$ 的阶为 $-1$，而 $P^0=Q^0=\mathbb Q_p$，其余为零，原右侧却为 $+1$。使用后面的首项归一化时也应同步留意此符号。}

\DeclareTranslatorNote{11}{Frobenius 同余条件的适用范围}{按所收勘误说明，同余条件应只施加于系数环：对 $a\in\mathcal O_K$ 要求 $\phi(a)\equiv a^q\pmod\pi$，再令 $\phi$ 固定形式参数，从而扩张到 $\mathcal O_K[[t]]$。原文同时要求固定 $\mathcal O(\Gamma)$，并对所有函数 $f$ 满足 $\phi(f)\equiv f^q\pmod\pi$；取形式参数 $t\in\mathcal O(\Gamma)$，便同时得到 $\phi(t)=t$ 和 $t\equiv t^q\pmod\pi$，在剩余域上的形式幂级数环中矛盾。后文使用的 $\phi^{-n}$ 应在修正后的系数 Frobenius 约定下理解。}

\DeclareTranslatorNote{12}{微分迹与拉回复合的标量}{此复合应为乘以 $[\pi]$ 的次数，即
\[
\Tr_\pi\circ[\pi]^*=q\,\id_{\Omega^1(G)},
\qquad q=\#(\mathcal O_F/\pi),
\]
而非乘以 $\pi$。这与 2.3.1 后所定义的迹及投影公式相符；在相邻层 $m=n+1$ 取 $x=1$ 即得到标量 $q$。例如 $F/\mathbb Q_p$ 为非分歧二次扩张、$\pi=p$ 时，该标量是 $p^2$，不是 $p$。}

\DeclareTranslatorNote{13}{Kummer--$d\log$ 方块的符号}{按 2.2.9 对连接同态的约定，图中左方块实际反交换。为使它交换，可把左侧竖直箭头改为 $-d\log$；另一种表述是保留箭头而明确称其反交换。证明中令 $w=\log(\widetilde v^{p^m}h^{-1})$，已算得 $dw=-dh/h$，而 $\mu(\omega)$ 对 $dx=\omega$ 由余循环 $\sigma x-x$ 表示。因此证明最后一行的类应为 $-\mu_{\xi,n}(dh/h)$，不能省掉负号。此处的修正须与前后显式互反律中的符号一致。}

\DeclareTranslatorNote{14}{$\log\chi_G$ 的值域与整性}{一般应先把 $\log\chi_G$ 看作 $H^1(K,F)$ 中的元素；若需整系数形式，再乘一个足够可除的固定常数。分歧情形中，$\log:\mathcal O_F^\times\to\mathcal O_F$ 的整性断言不总成立。例如取 $F=K=\mathbb Q_3(\pi)$、$\pi^4=3$，并令 $v_\pi(\pi)=1$。则
\[
\log(1+\pi)=\sum_{j\ge1}(-1)^{j-1}\frac{\pi^j}{j}
\]
各项的赋值为 $j-4v_3(j)$，唯一最小值在 $j=3$ 处为 $-1$，故 $\log(1+\pi)\notin\mathcal O_F$。Lubin--Tate 特征可取到 $\mathcal O_F^\times$ 中这些值，因而需另行处理整性。这与先前译注指出的对数二次项正负号是两个不同的问题。}

\DeclareTranslatorNote{15}{de Rham 实现缺少滤过上标}{第二个同构应读为
\[
\Omega\otimes\mathbb Q_p\simeq D^0_{\mathrm{dR}}(\mathscr F),
\]
即取滤过的第 $0$ 部分。第一章 3.2.6 所定义的 $\Omega$ 是 $M^0_{\mathrm{dR}}$，第二章 2.1.3 中相应的 $\operatorname{coLie}$ 同构也以 $D^0_{\mathrm{dR}}$ 为源。在复乘情形，完整 de Rham 空间与该滤过部分的秩可能不同，不能省去上标。}

\DeclareTranslatorNote{16}{辅助映射 $a'$ 并非一般同构}{此处 $a'$ 上的 $\xrightarrow{\sim}$ 应改为普通箭头；经指定的系数特化后得到的 $a$ 才是所需同构。例如在有理数情形取 $r=1$、$L=\mathbb Q(i)$ 和奇二次特征，所写 $a'$ 对应实线性映射
\[
z\longmapsto(z\bmod\mathbb R,\bar z\bmod\mathbb R).
\]
定义域和值域的实维数均为 $2$，但其核为 $\mathbb R$、秩为 $1$。投影到所需的奇特征分量后，才得到一维同构。不能据原式把未经特化的 $a'$ 也当作同构。}

\DeclareTranslatorNote{17}{复乘局部计算应取 $\mathfrak p$ 分量}{这段局部计算应使用
\[
T_{\mathfrak p}E
=T_pE\otimes_{\mathcal O_K\otimes\mathbb Z_p}\mathcal O_{\mathfrak p},
\]
并把相应 Betti 格换成 $H_1(E(\mathbb C),\mathbb Z)\otimes_{\mathcal O_K}\mathcal O_{\mathfrak p}$。例如 $p$ 在复乘域中分裂时，$\mathcal O_{\mathfrak p}=\mathbb Z_p$，而完整的 $T_pE$ 的 $\mathbb Z_p$ 秩为 $2$，不可能以所写的一个元素作为 $\mathcal O_{\mathfrak p}$ 基；投影到 $\mathfrak p$ 分量后才有相应的秩一表述。}

\DeclareTranslatorNote{18}{复乘情形的系数代数}{这里应取
\[
A=K[\Gal(L/K)],
\]
而不是 $\mathbb Q[\Gal(L/K)]$。本段处于虚二次域情形，前文要求 $A$ 是 $K$ 的有限扩张的有限直积，并且使用 $K[\Gal(L/K)]$ 上的张量积。若 $\mathbb Q[G]$ 有所需的含幺 $K$ 代数结构，将其与增广商 $\mathbb Q[G]\to\mathbb Q$ 复合就会给出 $K\hookrightarrow\mathbb Q$，与 $K$ 为虚二次域矛盾。}

\DeclareTranslatorNote{19}{Siegel--Bombieri--Vaaler 定理的高度归一化}{定义 2.3(d) 沿用的是 Plücker 坐标向量的 sup 高度；定理 2.4 中不增加修正常数的版本，在无穷位须使用 Euclidean Plücker 范数。若保留当前 $h(A)$，可在定理右侧增加 $\tfrac12\log\binom NM$ 来吸收两种矩阵高度之差。例如 $k=\mathbb Q$、$A=(5,7,8)$ 时，当前 $h(A)=\log8$，但任何非零有理核向量化成本原整数向量后，最大坐标绝对值至少为 $3$：若各坐标绝对值至多为 $2$，由 $y\equiv z\pmod5$ 得 $y=z$、$x=-3y$，只能为零。因此任一核空间基的高度和至少为 $\log9>\log8$；$(3,-1,-1)$ 与 $(1,-3,2)$ 达到该下界。Euclidean Plücker 高度在此例中为 $\tfrac12\log138$，不会产生上述矛盾。}

\DeclareTranslatorNote{20}{Riemann--Roch 下界的一次项}{这里的 Riemann--Roch 公式还含有一次项：
\[
\chi(\mathcal O(dY))=\delta d^2-(g-1)(a_1+a_2)d+(g-1)^2.
\]
由 $h^2=0$ 只能直接推出 $h^0\ge\delta d^2-O(d)$；也可在 $d$ 充分大时使用 $h^0\ge(\delta/2)d^2$，不能省略低阶项而声称 $h^0\ge\delta d^2$。例如 $g=2,a_1=4,a_2=1,\delta=2$ 时，$Y=\Delta+3F_1$ 丰沛，对充分大且满足整除条件的 $d$，有 $h^0(\mathcal O(dY))=2d^2-5d+1<2d^2$。这一修正保留所需的二次增长，但须相应调整后续常数。}

\DeclareTranslatorNote{21}{单项式范围的严格端点}{两个范围应分别读为 $0\le b<N$ 和 $0\le b'<N$，上界不能取到 $N$。以一个因子说明：若取 $0\le a+b\le D$、$0\le b\le N$，当 $D$ 大时，单项式数为
\[
(N+1)(D+1)-N(N+1)/2,
\]
其一次项为 $(N+1)D$；目标空间 $H^0(C,\mathcal O(DNF))$ 的维数却为 $ND+1-g$。源空间维数已经大于目标，不能得到所声称的单射。改为严格不等式后，后面的线性方程构造也应使用相同的范围。}

\DeclareTranslatorNote{22}{范数比较所需的测度假设}{除书末附记已提到的 $X$ 紧致性外，还须对测度作非退化假设。为与给出的局部积分证明相配，可要求 $\nu$ 为光滑、处处正的体积测度。仅说“固定一个测度”不够：在紧致的 $X=\mathbb P^1$ 上取一点的 Dirac 测度 $\delta_P$，再取 $\mathcal O(1)$ 的一个在 $P$ 消失的非零截面 $s$，就有 $\|s\|_{L^2}=0$ 而 $\|s\|_{\sup}>0$，不可能有正常数 $c_{\mathscr L}$ 满足所写的不等式。}

\DeclareTranslatorNote{23}{Faltings 复形 (9.3) 末项的符号}{式 (9.3) 末项中的 $(s_i\operatorname{pr}_i-s_j\operatorname{pr}_j)^*L$ 应读为带加号的 $(s_i\operatorname{pr}_i+s_j\operatorname{pr}_j)^*L$。令
\[
D_\pm=\sum_{i<j}(s_i\operatorname{pr}_i\pm s_j\operatorname{pr}_j)^*L,
\qquad H=\sum_i s_i^2\operatorname{pr}_i^*L.
\]
则 $D_-+D_+=2(n-1)H$。源线丛为 $A=d(D_-+\delta H)$，中项为 $B=d(2n-2+\delta)H=A+dD_+$；乘子 $u_\alpha$ 属于 $\mathcal O(dD_+)$，故关系 $u_\alpha t_\beta-u_\beta t_\alpha$ 应落在 $\mathcal O(B+dD_+)$。再乘以 $dD_-$ 的截面才得到 (9.4)。这与 §11 中已另行记录的符号问题不是同一处。}

\DeclareTranslatorNote{24}{范数截面多次数中的总次数}{式 (13.10) 应读为
\[
e_i=d_i\prod_{j=1}^nN_j=d_iN,
\qquad N=\deg(\pi_1\times\cdots\times\pi_n).
\]
这里的原线丛是 $\pi^*M$，$M=\mathcal O(d_1,\ldots,d_n)$，由 $\Nm_\pi(\pi^*M)=M^N$，每个多次数都应乘总次数 $N$，不能漏掉 $N_i$。仅取 $n=1$，原式给出 $e_1=d_1$，而范数实际给出 $e_1=N_1d_1$。推论 18.3 证明中再次写出的 $e_i=d_iN/N_i$ 也应同步按 $d_iN$ 理解。}

\DeclareTranslatorNote{25}{命题 14.1 中 $\Gamma_1$ 的 $\epsilon$ 符号}{这里应为
\[
\Gamma_1=\Gamma\!\left(\mathscr W,d(2n-2-\epsilon)\sum_i s_i^2\operatorname{pr}_i^*L\right)^a.
\]
所用复形来自把 $\delta=-\epsilon$ 代入 (9.4)，所以 $2n-2+\delta$ 应成为 $2n-2-\epsilon$。保留加号会使构造属于另一条线丛，无法给出命题所需的 $L_{-\epsilon,s}$ 的截面。}

\DeclareTranslatorNote{26}{非阿基米德 Hensel 引理的完备性}{此处应增加 $K$ 关于给定非阿基米德绝对值完备的假设，才可从 Newton 序列推出它在 $K$ 中收敛到根。例如 $K=\mathbb Q$ 配 $7$ 进绝对值，取 $f(T)=T^2-2$、$\alpha_0=3$，则 $|f(3)|_7=1/7<1=|f'(3)|_7^2$，满足所写条件，但根 $\sqrt2$ 不在 $\mathbb Q$ 中。后文实际使用的完备域 $\mathbb C_v$ 满足补充后的假设。}

\DeclareTranslatorNote{27}{复数 Newton 引理应假设全纯}{此处应要求各 $f_i$ 全纯。若只要求实意义下的 $C^2$，不能仅由 $m\times m$ 复 Jacobian 和 $\partial^2/\partial z_i\partial z_j$ 的估计推出结论。例如 $m=1$、$f(z)=z+\bar z+i/4$、$\alpha_0=0$、$B=1$ 时，$f_z=1$、$f_{zz}=0$，满足所写估计，但 $\operatorname{Im}f(z)=1/4$，根不存在。若要保留一般 $C^2$ 映射，应改用完整的 $2m\times2m$ 实 Jacobian 及相应 Hessian，并重新表述估计。}

\DeclareTranslatorNote{28}{乘积定理高度界所需的整性}{定理 18.1 及推论 18.3 的 (18.2)、(18.4)，只用阿基米德范数控制高度时，须补上数域及整性条件：截面应来自固定整数模型，且所用模型、格及度量须在估计中固定。另一种修正是把右侧改为计入有限位分母贡献的适当射影高度。否则把任意截面 $\gamma$ 换成 $\gamma/M$，零集、指标和次数不变，阿基米德对数范数却可趋于负无穷。例如在定理 18.1 中取 $n=1$、$\Pi=\mathbb P^1_{\mathbb Q}$、$e_1=1$、$\gamma'=X_0/M$、$x=[0:1]$，并令 $0<\epsilon_3\le1$；所需真子簇只能是固定点 $\{x\}$，高度固定，但 (18.2) 右侧随 $M\to\infty$ 趋于负无穷。因此不能对任意有理截面直接使用现有的高度界。}

\DeclareTranslatorNote{29}{范数估计 (18.5) 的总次数幂}{若有限映射总次数为 $N$ 且 $\gamma'=\Nm(\gamma)$，右侧第一个因子应为
\[
\left(\prod_{v\mid\infty}\|\gamma\|_{\sup,v}\right)^N.
\]
因此 (18.5) 应在该因子后保留原有的指数误差项。缩放可见这一幂次不可省略：把 $\gamma$ 换成 $m\gamma$，范数截面变为 $m^N\gamma'$；原左侧按 $m^N$ 增长，原右侧却只按 $m$ 增长，$N>1$ 时矛盾。后文取对数并使用 (18.5) 的估计也须同步计入这个 $N$；由于相关次数有界，吸收常数时仍须明确这一来源。}
